\documentclass[10pt,twocolumn]{ltjsarticle}
\usepackage[a4paper,margin=14mm,columnsep=7mm]{geometry}
\usepackage{amsmath,amssymb,amsthm,booktabs}
\usepackage[hidelinks]{hyperref}
\usepackage{url}
\newtheorem{theorem}{定理}
\newtheorem{lemma}{補題}
\newcommand{\send}{\mathsf{Send}}
\newcommand{\other}{\mathsf{Other}}
\newcommand{\tick}{\mathsf{Tick}}
\title{レート制限付きAPIクライアントの受動的計装の形式化\\
--- Isabelle/HOLによる事例研究 ---}
\author{著者情報は後日記入}
\date{草稿 --- 実験結果は未記入}
\begin{document}
\maketitle
\begin{abstract}
並行して動作するAPIクライアントでは、要求の合計レートを制御する必要が生じる。計装によってその振る舞いに関する証拠を得られるが、モニタ自体が制御の流れやタイミングを変える場合があり、観測の欠落は結論を無効にしうる。本稿では、DeepL APIを利用するSlack翻訳アプリを題材として、Isabelle/HOLによる事例研究を行う。離散的なトークンバケット遷移系と受動的なトレースモニタの関係を、機械検査された定理として示す。別々のモデルで、クライアントの交錯実行、回数制限付き再試行、許可から要求までの遅延を記述する。抽象レベルでは、観測器を消去しても再試行のトレースが保存され、適切に初期化されたモニタはモデル中のバケットと一致する。さらに、時刻区間と実行部分列の対応を明示的に仮定することで、遅延を考慮した要求数の上限を条件付きで導く。独立した実験用アプリには３種類の観測モードを実装し、故障の隔離をローカルテストで確認するが、TypeScriptの精密化は証明していない。SlackとDeepLを用いる暫定的な評価計画を示し、実測結果は今後追記する。提案するレート方針は研究上の仮定であり、DeepL内部のレート制限に関する主張ではない。
\end{abstract}
\noindent\textbf{キーワード：} Isabelle/HOL、実行時検証、トークンバケット、APIクライアント、計装、Slack

\section{はじめに}
現代のアプリケーションでは、独立してスケジュールされる複数のタスクから外部APIを呼び出すことが多い。トークンバケットのような方針によって要求の合計レートを制限できるが、実装では再試行、タイミング、観測記録の欠落にも対処しなければならない。呼び出しを記録する観測器自体もソフトウェアであり、時計、識別子生成器、ログ出力先が故障する可能性がある。また、計測の追加は実行時間を変えうる。したがって、レート方針のみをモデル化しても、計装後のアプリケーションが計装前と同じ振る舞いを持つとはいえない。

本研究では、DeepLbot3に基づくSlack翻訳ボットを事例として、この違いを検討する。リアクションによってワークフローが起動し、DeepL翻訳APIを呼び出して返信を投稿する。既存ボットの変更を避けるため、独立した実験用コピーDeepLbot3xを使用する。クライアントは\texttt{fetch}を呼び出す。本稿の抽象的な要求イベントはこの呼び出しを表し、遠隔サーバへの到着を表すものではない。モデル中の方針とパラメータは研究のために設定する。DeepLは利用上限等を公表しているが~\cite{deepllimits}、本稿は公開されていないサービス側の要求レート方針を推定しない。

本稿の貢献は次の４点である。（１）レートバケット、交錯実行、再試行、受動的観測を対象とする有限遷移系とIsabelle/HOLの証明をまとめる。（２）共有バケットの実行と、遅延のあるfetch回数の上限を条件付きで接続する。（３）独立した実験用ボットについて、モデルと実装の対応付け、および故障を対象としたテストを示す。（４）その対応と計測負荷を実環境で評価する暫定計画を示す。これらの定理は、TypeScript実装が形式モデルを精密化すること、Slackログが完全であること、提案した分散レート方針がデプロイ済みアプリで強制されることを証明するものではない。

\section{背景と関連研究}
ネットワーク計算は、トラフィック制約と決定論的な上限を研究する~\cite{network,regulators}。本稿のトークンバケットはこの考え方に沿うが、方針を仮定として扱う離散的な操作的モデルである。本稿が証明するのはモニタとその操作的モデルとの等価性であり、独立に定義された到着曲線仕様との等価性ではない。

実行時検証は、実行の観測結果を仕様に照らして解析する~\cite{runtime}。本研究は、抽象的な受動観測器と観測対象の遷移系との関係、および実装から得られた観測をモデルに基づいて解釈するための条件に焦点を当てる。記録がないことを、イベントが発生しなかったことと同一視しない。

線形化可能性は、並行操作を適切な順序で原子的に実行されたかのように考えるための枠組みを提供する~\cite{linear}。本稿の並行モデルは共有バケットへの原子的アクセスを仮定するが、分散リミッタを実装したり、その線形化可能性を証明したりするものではない。検証済み時間オートマトンモデル検査~\cite{timed}はタイミングの仮定に注意を払う動機を与えるが、本稿は時間オートマトンのモデル検査器ではなく、手で定義した離散時間遷移系を用いる。証明環境にはIsabelle/HOL~\cite{isabelle}を用いる。本稿の貢献は、一般的なトラフィック制御や実行時検証の新アルゴリズムではなく、対象アプリに即した統合と仮定の明示である。

\section{システムとモデル化の範囲}
対象の元ボットはSlackのDeno SDKを用いたTypeScriptアプリケーションである。リアクションがワークフローを起動し、メッセージの情報を取得し、翻訳クライアントを呼び出し、返信を投稿し、その後の記録処理を行う。Slackのカスタム関数は、アプリケーションが定義した処理をワークフローのステップとして実行する~\cite{slackfunctions}。調査対象のクライアントは翻訳の試行ごとに１回fetchを呼び出す。以下の共有リミッタと再試行のプロトコルは将来の実装を想定したモデルであり、元ボットに存在する機能ではない。

本稿は、（１）抽象遷移系、（２）モデル中のイベントをトークンバケットのトレースに写す観測関数、（３）fetch境界に手書きで実装したTypeScript観測器、（４）Slackによるログ配送、の４段階を区別する。形式的な定理は（１）と（２）に適用される。単体テストは（３）の限定的な証拠を与える。証明もテストも（４）の完全性を確立しない。fetch観測の成功はローカルな呼び出しを表し、サーバ側での受信、処理完了、課金を意味しない。

\section{トークンバケットと実行時モニタ}
容量を$B$、補充率を$r$、残量を$b$、時間増分を自然数とする。$b\leq B$の下で、時間イベント$\tick(d)$は残量を$\min(B,b+rd)$に更新する。送信イベント$\send$は$b>0$のときのみ可能で、残量を１減らす。$\other$は残量を変えない。有限イベント列$E$中の送信回数を$S(E)$、時間増分の総和を$T(E)$とする。有効な実行を$b\xrightarrow{E}b'$と記す。

\begin{lemma}[資源収支]\label{lem:account}
$b\xrightarrow{E}b'$ならば、$S(E)+b'\leq b+rT(E)$が成り立つ。
\end{lemma}
\begin{proof}[証明の考え方]
$E$について帰納法を用いる。送信は$S$を１増やし残量を１減らす。時間イベントによる残量の増加は高々$rd$である。その他のイベントは両者を変えない。
\end{proof}
\begin{theorem}[有限実行における上限]\label{thm:bucket}
$b\leq B$かつ$b\xrightarrow{E}b'$ならば、$S(E)\leq B+rT(E)$である。
\end{theorem}
\begin{proof}[証明の考え方]
補題~\ref{lem:account}と$b'$の非負性を用いる。
\end{proof}
これらの結果は\texttt{Rate\_Bucket.thy}で形式化されている。\texttt{Rate\_Monitor.thy}のモニタは、残量$\mathsf{Some}(b)$または永久的な棄却状態$\mathsf{None}$を保持する。同じ遷移に従って更新し、残量０で送信を観測すると棄却する。
\begin{theorem}[操作的モニタとの対応]\label{thm:monitor}
任意の有限列$E$について、$b\xrightarrow{E}b'$であることと、$\mathsf{Some}(b)$から開始したモニタが$E$の処理後に$\mathsf{Some}(b')$を返すことは同値である。
\end{theorem}
したがって、$b\leq B$から始まる受理トレースは定理~\ref{thm:bucket}を満たす。ただし、この結論は\emph{観測された}トレースに関するものであり、欠落したイベントについては何も示さない。

\section{並行実行、再試行、受動的観測}
\subsection{クライアントの交錯実行}
\texttt{Concurrent\_Bot.thy}では、実行識別子から処理段階への全関数により、潜在的に多数のクライアントを表す。各クライアントの局所動作は、１つの共有トークン残量の下で交錯する。抽象的な境界では、要求動作が原子的にトークンを消費する。射影定理によって並行実行は有効なバケット実行に写されるため、定理~\ref{thm:bucket}がモデル中の全送信数に適用される。動的生成や具体的な共有状態の実装は含まれない。

局所状態と共有状態の違いは重要である。容量１の独立したバケットが２つあり、補充も時間経過もない場合、各クライアントは１回ずつ送信できる。合計は２回となるが、１つの共有バケットなら１回しか許可されない。機械検査されたこの反例は、クライアントごとの検査だけでは全体のレート上限を確立できないことを示す。

\subsection{再試行と期限}
\texttt{Retry\_Bot.thy}は、ジョブ状態、試行回数、再試行可能時刻、期限を共有バケットとともに別個にモデル化する。送信には、待機中であること、トークンがあること、試行回数が$K$未満であること、および時刻が期限より厳密に前であることが必要である。レート制限応答を受けると、モデル中の２種類の待機時間の最大値を用いて次の可能時刻を設定する。クォータ超過、タイムアウト、成功は実行中のジョブを終了し、期限切れは待機中ジョブを終了する。証明済みの性質には試行回数の上限、送信条件、終了状態の吸収性が含まれる。タイムアウトはクライアント側の観測であり、サーバ側での取消しや課金の不存在を示さない。公平なスケジューリング、最終的な進行、翻訳成功は証明していない。このモデルが並行ボットモデルを精密化することも示していない。

\subsection{消去と対応}
\texttt{Instrumented\_Retry.thy}では、再試行の実行にモニタ状態を付加する。各再試行イベントをバケットイベントに射影し、基となる再試行遷移を変えずに観測器を更新する。再試行実行を$R(s,E,s')$、観測付き実行を$O(s,m,E,s',m')$、モニタを$M$とすると、機械検査された\texttt{observational\_correspondence}定理は次を述べる。
\begin{equation}\label{eq:erase}
 O(s,m,E,s',m')\iff R(s,E,s')\land m'=M(m,\pi(E)).
\end{equation}
観測器を消去しても$E$と再試行系の最終状態は変わらず、すべての再試行実行に観測付き実行を付加できる。モニタが棄却状態でもモデルの実行は妨げられない。抽象バケット残量に一致するよう初期化すれば、モニタの最終残量も一致する。イベントが正確に配送されること、単一の抽象時間軸、観測器が時間を消費せず例外を出さないことは構成上の前提である。式~\eqref{eq:erase}はTypeScriptの実行や物理的なタイミングに関する定理ではない。

\section{許可からfetchまでの遅延}
各fetchに一意の許可識別子が対応し、各許可が高々１回使われるとする。許可時刻を$g$、fetch時刻を$f$とし、$g\le f\le g+D$を仮定する。また、任意の半開区間$[s,t)$の許可回数が$B+r(t-s)$以下であるとする。使用されない許可があってもよい。
\begin{theorem}[遅延を考慮した上限]\label{thm:delay}
任意の$s\leq t$に対し、$[s,t)$内のfetch回数は$B+r((t-s)+D)$以下である。
\end{theorem}
\begin{proof}[証明の考え方]
各fetchは$[\max(0,s-D),t)$内の許可へ単射的に対応する。許可回数の上限を適用し、区間幅が高々$(t-s)+D$であることを用いる。
\end{proof}
機械検査された\texttt{Permit\_Delay.thy}は、遅延を考慮しない元の上限への反例も構成する。$B=r=1$で、時刻０、１、２の許可をすべて時刻２にfetchすると、$D=2$を満たしながら区間$[2,3)$で元の上限を超える。したがって、許可時点の上限を、後のfetch時点の同じ上限として自動的に解釈することはできない。

\texttt{Permit\_Operational.thy}は許可、fetch、取消しを別々にモデル化し、一意の識別子とfetch時の遅延条件を扱う。その不変条件は、許可が最終的に使われることや、実際のスケジューラが$D$を満たすことを確立しない。\texttt{Permit\_Bridge.thy}は並行実行の任意の連続部分列の上限を証明し、\texttt{window\_correspondence}の下で遅延定理に接続する。この仮定は、各許可時刻の区間に、許可数が一致し経過時間が区間幅以下の連続部分列が対応することを要求する。これは遷移系からまだ導出されていない。また、別々の実行が同じ実際の実行を表すという解釈も必要である。この対応の導出と、再試行モデル・許可モデルの接続は未解決である。

\section{実装との対応と故障の隔離}
実験用パッチは、固定されたDeepLbot3のリビジョンに基づく~\cite{repository}。注入されたfetch関数をラップして、ローカルな単調時計による時刻とメタデータを記録し、必要に応じて１回の関数実行内に限定された容量１・補充０の観測器を動かす。モニタは受動的であり、呼び出しを抑止しない。この具体的な観測器では、別々のSlack Functionにまたがる全体上限を確立できず、実行ごとに異なる時計の起点を相互に比較できるとも仮定できない。許可に関する理論に対応する許可操作は実装されていない。

３つの実験モードは共通の計測コードを持つ。baselineには識別子生成、時刻取得、文字数計数、共通ログ出力が含まれ、recordはfetch観測を追加し、monitorは局所モニタの判定を追加する。したがってbaselineとの比較は、\emph{共通計装を含む}対照に対する差を測るものであり、無計装プログラムとの総負荷差ではない。観測器の計算時間はログ出力を含まず、ログ処理全体の負荷ではない。

故障注入によって、時計、識別子生成、文字数計算、ログのシリアライズ、エラー分類が失敗すると振る舞いが変わりうることが判明した。修正後の候補では同期的な観測操作を保護し、取得できない値を０と偽ることなく欠測として扱う。テストは一部の時計・出力先の故障、特殊な拒否値、およびfetchの元の結果やエラーの伝播を対象とする。これは選択したケースのテストであり、すべての実行に対する証明ではない。プロセス停止、メモリ枯渇、非同期ログ故障、配送損失は実証範囲外である。研究リポジトリにはオフラインのパッチ検証と候補コードの来歴を記録している~\cite{repository}。

\section{実験方法と暫定的な状況}\label{sec:experiment}
\textbf{暫定稿：}実証研究は進行中である。以下の実験計画と数値は予定であり、完了した結果ではない。旧版の実験パッチについて以前行った動作確認では、異なる入力を用いて各モード１回の翻訳が成功した。この観測は修正後の候補コードを検証するものでも、定量的な比較を支えるものでもない。

\subsection{研究課題と環境}
３モードの実行成功と失敗の振る舞い（RQ1）、記録されたイベントを試行識別子に対応付けられるか（RQ2）、計測された処理時間にどのような差があるか（RQ3）を調べる予定である。独立してインストールした実験用Slackアプリ、専用テストチャンネル、合成日本語テキスト、DeepL API Freeの接続先を実環境として用いる。ソースリビジョン、パッチの識別情報、SDKのバージョン、実行環境、設定、日付、計画からの逸脱を記録する。認証情報、入力・翻訳文、ワークスペースやユーザーの識別情報は公開用の測定データに含めない。

本実験の前に、各モード１回の動作確認で統合とログ抽出を調べる。実行間隔、設定反映、ログ収集、復旧方法を決めるため、最大30回の予備試行を計画する。本実験の候補は約３日間に12セッションを設け、各セッションで各モード10回、合計360回を行うものである。セッション間でモード順序を均衡させ、各モードに同じ合成入力集合を用いる。動作確認３回、予備試行30回、本実験360回を合わせると、再実行を除き最大393回となる。開始間隔15秒は暫定案であり、予備試行による検証が必要である。環境変数によるモードと試行識別子の変更には反映確認が必要であり、実行の重なりや設定との対応付けに問題があれば本実験前に設計を改める。

\subsection{測定と分析の計画}
予定された試行と観測された試行、モード、匿名識別子、設定、開始・終了のメタデータ、結果、エラー、欠測、取得可能な関数・翻訳・観測器の時間を記録する。可能な場合にはアカウント全体の使用量を別に確認する。送信したコードポイント数と課金文字数が等しいとは主張しない。試行予定の台帳や連番は一部の欠落を発見できるが、すべてのfetch呼び出しが観測されたことや、最後のイベントが配送されたことを証明できない。欠測値はnullのままとし、完全に失われた実行を観測成功として数えない。

各モードの標本数、成功・失敗・欠測、および遅延分布とセッションごとの変動を示す。比較では入力を対応付け、順序とセッション内の相関を考慮する。高いパーセンタイルの分析では標本数の限界を明記する。Function内の経過時間はSlack全体の応答時間とは限らない。例外隔離に関する限定的なテストは、すべての環境での意味保存を示さない。予備試行と動作確認は本実験のデータと分ける。本実験の前に停止条件、再試行、再実行、除外条件を定めて記録する。

\subsection{実験結果の記入欄}
\emph{実験の実施と計画の検証後に追記する。}本草稿では、実環境の性能、観測の網羅率、統計的有意性、全体的な並行制御について数値による主張を行わない。

\section{限界と今後の課題}
抽象モニタは、その構成元であるバケット意味論に対応する。形式化は、独立した仕様、分散環境での原子性、実装の精密化、物理的な時間保存、イベント配送を検証していない。遅延に関する接続は未証明の区間対応を仮定しており、再試行モデルと許可モデルは統合されていない。実験用観測器は実行ごとに存在し、クライアントには共有レートリミッタも再試行ロジックもない。実験は実装の欠陥を発見し、試験した条件を特徴付けられるが、普遍的な性質を証明したり、不明なサービス提供者の方針を推定したりすることはできない。今後は、共通の時刻付きトレースから接続条件を導き、欠落を考慮した観測を仕様化し、実際の共有リミッタとの実装対応を証明する必要がある。

\section{おわりに}
本草稿は、レート制限付きAPIクライアントの受動的観測に対する形式的・実験的アプローチを示した。Isabelle/HOLの結果は、共有された抽象バケット、正確なイベント配送、明示的な遅延の仮定から何が従うかを示す。実験計画は、ログ欠落、観測負荷、未証明の接続を明らかにしたまま具体的な実装を調べる。実環境での評価とその解釈は今後の課題である。

\small
\bibliographystyle{plain}
\bibliography{references}
\end{document}